\scanpage{115}
\mspage{57}

s'identifie bel et bien au groupe de Galois $\operatorname{Gal}(\overline K/K)$ de \deleted{$\unsure{\mathcal E_Y}$} $\overline K/K$, $\overline K$ étant la clôture sép. de $K$ telle que $\operatorname{Spec}\unsure{\widetilde K|\overline K}=\widetilde Y$. Souvent, on notera $\Gamma$, \inserted{ou $\Gamma_K$, $\Gamma_K^{\overline K}$,} au lieu de $\mathcal E_Y$ --- surtout si $K$ est alg. sur le corps premier, et qu'il n'a donc plus de ``partie géométrique'' à distinguer d'une ``partie arithmétique''$\ldots$

Soit \inserted{\unsure{un}} $K$ \deleted{\unsure{tel}} corps (qui pourrait être alg. clos), $L$ une ext. de type fini de $K$, $X$ un ``modèle'' propre régulier de $L$. Alors $\mathcal E_X^{\overline L}$ s'identifie à un quotient de $\mathcal E_L^{\overline L}$, qui \uline{ne \deleted{\unsure{p}} dépend pas du modèle $X$ choisi}, comme il est bien connu

\margin{114}
